\section{Conventions and birational inputs}
\label{sec:preliminaries}

Varieties are integral and projective over an algebraically closed
field of characteristic zero unless stated otherwise. A $\Q$-pair
$(X,B)$ consists of a normal variety, an effective rational boundary
$B$, and a rational Cartier divisor $K_X+B$. We use log discrepancies:
on a resolution $h:Y\to X$, the equality
$K_Y+B_Y=h^*(K_X+B)$ assigns log discrepancy $1-b_E$ to a prime divisor
with coefficient $b_E$ in $B_Y$. The pair is klt if all log
discrepancies are positive, and log canonical if they are nonnegative.
We use the standard cone, contraction, negativity, and resolution
theorems as in \cite{KM,BCHM}.

The symbols $\sim_\Q$ and $\equiv$ denote rational linear and numerical
equivalence. A divisor is pseudo-effective if its numerical class is
in the closure of the cone generated by effective divisors, and big if
its class lies in the interior of that cone. Its Iitaka dimension is
denoted by $\kappa$. Over $\C$, a very general point means a point
outside a specified countable union of proper closed subsets.

For a nef divisor $N$, its nef reduction is an almost holomorphic
dominant rational map with connected compact general fibres, on which
$N$ is numerically trivial, such that every curve through a very
general point not contracted by the map has positive $N$-degree
\cite{NefReduction}. Its base dimension is the \emph{nef dimension}
$n(X,N)$. Thus $n(X,N)=\dim X$ precisely when every curve through a
very general point has positive $N$-degree. This numerical condition
does not assert bigness.

For a pseudo-effective real divisor $D$ on a smooth variety and a
prime divisor $\Gamma$, fix an ample divisor $A$ and set
\[
 \sigma_\Gamma(D)=\lim_{\delta\downarrow0}
 \inf\{\mult_\Gamma D':D'\ge0,\ D'\sim_\R D+\delta A\}.
\]
These numbers are independent of $A$ and depend only on the numerical
class of $D$. The divisorial Zariski decomposition is
\[
 N_\sigma(D)=\sum_\Gamma\sigma_\Gamma(D)\Gamma,
 \qquad P_\sigma(D)=D-N_\sigma(D).
\]
The sum has finite support. A nef divisor has $N_\sigma(D)=0$.
Further properties from \cite{Nakayama,LP} are recorded at their point
of use in Section~\ref{sec:positive}.

\begin{lemma}\label{lemma:pseudo-vg}
Let $D$ be a pseudo-effective rational Cartier divisor on a projective
complex variety $X$. Every integral curve through a sufficiently very
general point has nonnegative $D$-degree.
\end{lemma}
\begin{proof}
Fix an ample Cartier divisor $A$ and positive rational numbers
$\delta_i\to0$. Each $D+\delta_iA$ is big, so there is an effective
rational divisor $E_i\sim_\Q D+\delta_iA$. Choose a point outside
$\bigcup_i\Supp E_i$. A curve through it is contained in none of
these supports; hence $(D+\delta_iA)\cdot C\ge0$ for every $i$.
Let $i\to\infty$.
\end{proof}

We state explicitly the companion results used throughout. We require
existence of a terminating choice of program, rather than termination
of every sequence of flips.

\begin{theorem}[Log abundance, {\cite[Theorem~1.1]{LA}}]
\label{input:abundance}
Let $(X,B)$ be a projective log canonical $\Q$-pair over an
algebraically closed field of characteristic zero. If $K_X+B$ is nef,
then it is semiample.
\end{theorem}

For the next input a nef rational Cartier b-divisor is specified by a
nef rational Cartier divisor on a projective birational model; it is
pulled back unchanged on all higher models. Its trace on a lower model
is its pushforward and need not be nef. Generalized discrepancies are
defined by pulling back the adjoint while keeping this fixed nef data.

\begin{theorem}[Minimal-model programs,
{\cite[Theorem~1.1 and Section~3.3]{MM}}]
\label{input:models}
Let $(X,B+\mathbf M)$ be a projective generalized klt $\Q$-pair over $\C$,
with $X$ $\Q$-factorial and with fixed nef rational Cartier b-divisor
data. If $K_X+B+M_X$ is pseudo-effective, there is a terminating choice
of its minimal-model program, with $\Q$-factorial generalized klt
models and a nef endpoint. The program extracts no divisors and
generalized log discrepancies do not decrease. On a common resolution
of its initial and final models, the pullback of the initial adjoint
is the pullback of the final adjoint plus an effective divisor
exceptional over the final model.
\end{theorem}

The ordinary klt case follows by taking the nef data to be zero. In
that case Theorem~\ref{input:abundance} makes the endpoint a good
minimal model. The moving-center and Frobenius results used from
\cite{LA} are stated separately in
Sections~\ref{sec:jet-tools} and~\ref{sec:frobenius}; neither theorem
above asserts abundance for a nonzero generalized nef summand.
